\documentclass[11pt,a4paper]{article}
\usepackage[margin=24mm]{geometry}
\usepackage[T1]{fontenc}
\usepackage{lmodern,amsmath,amssymb,booktabs,longtable,array,microtype}
\usepackage[hidelinks]{hyperref}
\urlstyle{same}
\setlength{\emergencystretch}{2em}
\title{Prior Work and Its Connection to Our HNC/MD Study\\
\large Entropy, diffusion, and structural diversity in charged mixtures}
\author{EntropicFluids collaboration}
\date{31 October 2026}
\begin{document}
\maketitle
\begin{abstract}
This review collects the prior work discussed in connection with the
EntropicFluids study: excess-entropy scaling in simple and charged mixtures,
high-entropy electrolytes, and Voronoi/graph estimates of structural entropy.
It distinguishes what these studies establish from what remains to be tested
in our model mixtures. The general idea that multicomponent electrolytes can
show enhanced ion mobility already has direct precedent. The potential value
of the present study lies in its controlled comparisons, independent
thermodynamic-integration benchmark, and distinction between species labels,
interaction diversity, and probe-ion motion.
\end{abstract}

\section{Scope and terminology}
This is an annotated literature assessment rather than a systematic review
or a claim of priority. References include the studies discussed previously
and supporting methodological sources. The companion
\path{voronoi_entropy_work_plan.tex} gives the proposed analysis protocol;
this document supplies the scientific context.

Three entropy notions must remain distinct. Ideal compositional mixing has
entropy per particle $s_{\rm mix}^{\rm id}=-k_B\sum_i x_i\ln x_i$.
Thermodynamic excess entropy is
\begin{equation}
s_{\rm ex}=\frac{S-S_{\rm id,mix}(T,\rho,\{x_i\})}{N},
\end{equation}
so that the ideal mixing term cancels against the same-composition reference.
Finally, a graph or local-environment entropy describes the diversity of a
chosen structural representation. A solvation-state entropy is likewise
conditional on the chosen definition of a solvation state. These quantities
can correlate without being interchangeable. In the glass literature,
configurational entropy often refers specifically to basin multiplicity
after separating vibrational entropy, which is another distinct target.

\section{Excess entropy and diffusion in mixtures}
\subsection{Rosenfeld: foundations and charged-mixture precedent}
Rosenfeld's original work~\cite{Rosenfeld1977} motivates an approximate
relation between reduced transport and excess entropy. For self-diffusion,
write
\begin{equation}
D_i^{\rm R}=D_i\rho^{1/3}\sqrt{m_i/(k_BT)},\qquad
\ln D_i^{\rm R}\simeq a_i+b_i s_{\rm ex}/k_B.
\end{equation}
This representation is useful empirically; it does not imply that changing
any contribution called entropy must accelerate every species.

Rosenfeld, Nardi, and Zinamon~\cite{Rosenfeld1995} developed a corresponding-states
hard-sphere description of diffusion in binary dense-plasma mixtures, tested
against molecular dynamics of asymmetric ionic mixtures. Charged-mixture
entropy--diffusion connections therefore have a long history. Their plasma
setting differs from our explicit neutral solvent with oppositely charged
probe ions, but a broad claim of introducing entropy scaling to charged
mixtures would be inappropriate.

Dyre's perspective~\cite{Dyre2018} explains the connection with hidden scale
invariance and isomorph theory, while reviewing failures of universality.
Its practical implication here is that a successful fit over a restricted
composition range is evidence of an empirical relation, not an exact law
covering changes in Hamiltonian or composition.

\subsection{Individual-species diffusion: Krekelberg et al. (2009)}
Krekelberg et al.~\cite{Krekelberg2009} studied tracer diffusivities in
hard-sphere and Gaussian-core model fluids, including mixtures with
size asymmetry and nonadditive interactions. Their generalized Rosenfeld
reduction uses interaction information and dilute-limit thermodynamic and
transport behavior. It captures many composition effects, but encounters
limitations when the component diffusivities decouple strongly.

This is the main precedent for our distinction between the composition-weighted
mean self-diffusion and divalent-ion self-diffusion. Ordinary Rosenfeld units
and the generalized reduction are different constructions. Our measurements
should not be described as a test of their generalized prediction unless
its required dilute-limit quantities are actually calculated. A weak
species-level correlation with one bulk entropy is not intrinsically
inconsistent with mixture scaling research.

\subsection{A predictive mixture framework: Schmitt et al. (2025)}
Schmitt, Hasse, and Stephan~\cite{Schmitt2025} combine entropy scaling,
limiting pure-component and infinite-dilution transport properties, and
mixing rules to describe self-diffusion and mutual diffusion in binary
mixtures. Mixture entropy is supplied by an equation of state. They test
model and real-fluid systems, and emphasize that a general single-variable
mixture dependence does not persist.

Their work shows how entropy can remain useful even when it is insufficient
alone. It also reinforces a necessary distinction: our
$D_{\rm net}=\sum_i x_iD_i$ is a weighted mean of self-diffusion coefficients,
not a mutual-diffusion coefficient or an electrical conductivity. A direct
comparison with their predictive framework would require additional mixture
transport information and an extension beyond the binary demonstrations.

\section{High-entropy electrolytes: the closest application precedents}
\subsection{Wang et al.: multisalts and enhanced lithium mobility (2023)}
Wang et al.~\cite{Wang2023} compare a four-salt lithium electrolyte with
single-salt counterparts at matched total salt concentration. Experiments
and MD support altered solvation and enhanced lithium mobility, alongside
improved battery performance. The paper explicitly recognizes the challenge
of separating the effect of entropy from the properties of the selected
chemical components.

This closely matches the motivation of our simulations. However, chemical
substitution changes interactions as well as compositional diversity. The
reported association is therefore not, by itself, evidence that ideal mixing
entropy independently causes a mobility increase. Our identical-interaction
label control and controlled distribution of solvent interactions address
that attribution problem in a simpler setting.

A related paper by Wang et al.~\cite{WangAdvanced2023} examines entropy-driven
multisalt electrolytes, including solubility, solvation, interphase stability,
and lower-temperature operation. It broadens the application context but
is a separate paper from Ref.~\cite{Wang2023}; their findings should not be
combined into a single bibliographic attribution.

\subsection{He et al.: assessment from the supplied full article (2026)}
\label{sec:HeAssessment}
The supplied ten-page article~\cite{He2026} permits a more specific assessment
than the earlier abstract-level discussion. The principal evidence is in
Sec.~2.1 (journal pp.~683--685), Sec.~2.2 (pp.~685--686), and Fig.~3.
The supplementary information was not supplied or successfully retrieved;
claims below about reporting and uncertainty refer to the main article,
not to an assertion that the supplementary methods contain no additional checks.

\subsubsection{Design and strengths}
The authors compare four formulations, EL1, EL2, EL3, and EL5, with one,
two, three, and five anion species, respectively. Total lithium-salt
concentration remains 1.0~M, with the same carbonate solvent base and 5\%
FEC additive. EL1 contains 1.0~M LiPF$_6$; EL2 has 0.6~M LiPF$_6$ and
0.4~M LiTFSI; EL3 has 0.6~M LiPF$_6$, 0.2~M LiTFSI, and 0.2~M LiDFOB;
EL5 has 0.6~M LiPF$_6$ plus 0.1~M each of LiTFSI, LiDFOB, LiFSI, and
LiNO$_3$. Matching total salt concentration is a useful control. With
monovalent ions it also matches nominal ionic strength under complete
dissociation; actual association need not match.

The combination of Raman spectra, MD coordination statistics, representative
cluster calculations, and electrochemistry supports an altered solvation
environment. The solvent contribution to Li coordination decreases and
several anions participate directly. The application results on nitrate
solubility and electrode interphases add practical value. These strengths
should be assessed separately from the stronger mechanistic claim that
entropy determines diffusion.

Changing salt identities and fractions also changes binding, geometry,
association, and solvent exchange. Thus the series is a meaningful chemical
design study, but it does not isolate component number or an entropy change
as an independent intervention. Fixed molarity alone does not fix particle
number density, force-field interactions, viscosity, or molecular packing.

\subsubsection{The two entropy measures have different targets}
Their Eq.~(2) defines a solvation-state Shannon entropy,
\begin{equation}
S_{\rm sc}=-R\sum_a p_a\ln p_a,
\end{equation}
which rises from 32.5 to 40.4~J~mol$^{-1}$~K$^{-1}$ between EL1 and EL5.
This usefully summarizes the diversity of the selected Li-centered solvation
classes. Its value depends on the class definitions, neighbor cutoffs,
chemical labels, and sampling of rare classes. Adding anion identities
expands the set of possible chemical classes. A composition-preserving
random-label reference, or a comparison using merged chemical classes,
would help separate this combinatorial diversity from changed local ordering.
That suggestion does not imply that the real anions are mechanically identical.

$S_{\rm sc}$ is not the total liquid entropy: it omits within-class information
and correlations between overlapping Li environments. Its natural reference
is a mole of sampled Li environments, whereas our TI entropy is normalized
per particle of the entire mixture. Their Eq.~(1), the ideal mixing formula,
is likewise not a general expression for full configurational entropy in
an interacting molecular liquid. None of these quantities should be added
together without a consistent decomposition and normalization.

For the second measure, the prose on p.~686 specifies Li--anion and
Li--solvent pair correlations in the local solvation environment. Their
Eq.~(4), however, has the form of the standard all-species pair-entropy sum:
\begin{equation}
\frac{s_2}{k_B}=-2\pi\rho\sum_{\alpha\beta}x_\alpha x_\beta
\int_0^\infty\!\left[g_{\alpha\beta}(r)\ln g_{\alpha\beta}(r)
-g_{\alpha\beta}(r)+1\right]r^2\,dr.
\end{equation}
Visual inspection confirms that the printed bracket is
$g\ln g-[g-1]$: its sign is correct. The issue is the estimator's scope,
not a sign error in this equation.

If only the specified Li-containing pairs are retained, the result is a
selected contribution to $s_2$, rather than the total mixture $s_2$ or
$s_{\rm ex}$. If the implemented calculation actually includes all species
pairs, the prose needs clarification. The main article does not resolve
this ambiguity sufficiently for reproduction. The atom-versus-molecule
species convention, total density, mole fractions, pair multiplicities,
radial cutoff, and per-Li versus per-total-particle normalization must also
be established before comparing values across changing anion chemistries.
For molecular models these conventions matter even at the same salt molarity.

The assertion that the pair term accounts for more than 90\% of excess
entropy is not independently benchmarked against a total-entropy calculation
for these electrolytes in the main article. Such a fraction is not a general
identity. It is particularly insufficient to equate a selected local pair
sum with full excess entropy. A TI or other independently validated
thermodynamic reference is needed to measure the missing contribution.

\subsubsection{What the reported transport correlation establishes}
The main text reports an increase in Li self-diffusion from
$1.42\times10^{-7}$ to $2.38\times10^{-7}$~cm$^2$~s$^{-1}$, approximately
68\%, while the stated entropy proxy changes from approximately $-0.071$
to $-0.068$ in units of $S_2/(Nk_B)$. Fig.~3d plots $\log_{10}D_{\rm Li}$
against this pair proxy for four formulations. It is a semilog correlation,
not a straight-line test of untransformed $D$. Four means and a fitted
slope/intercept provide only two residual degrees of freedom; independent
replica uncertainty, fit confidence intervals, and held-out predictions are
not displayed in that figure. This observation does not establish whether
such checks were performed elsewhere.

Their Eq.~(5) uses dimensional $D$, whereas conventional Rosenfeld scaling
uses $D^{\rm R}=D\rho^{1/3}\sqrt{m/(k_BT)}$ and an intensive excess entropy.
At identical temperature, Li mass, and particle number density, ratios of
$D$ and $D^{\rm R}$ would coincide. Fixed salt molarity does not alone
establish the last condition for these changing molecular compositions.
A common, documented density convention and explicit normalization of the
entropy coordinate are therefore needed. This limitation does not mean that
the plotted trend is spurious; it limits its interpretation as a classical
Rosenfeld validation or a comparison with our bulk-entropy slopes.

The data support a correlation between formulation, selected structural
statistics, and Li mobility. They do not independently determine whether
entropy, altered binding, viscosity, association, or exchange kinetics is
the causal factor. Improvements in interfacial resistance and battery
cycling support the electrolyte design, but also depend on changed interphase
chemistry and are not independent validations of a bulk self-diffusion law.

\subsubsection{Relation to our HNC/MD work and useful follow-up}
The closest comparison is with an ion-centered pair or neighborhood
\emph{descriptor}, rather than with our total TI entropy. Their selected
local coordinate may be useful even when bulk entropy predicts probe
mobility imperfectly. It also motivates the proposed charge/species-aware
Voronoi analysis, with a chemical-composition reference.

Our independent TI benchmark and exact label null address attribution and
normalization questions that the main article leaves open. Our controlled
interaction-amplitude comparison is also more restrictive, although it still
changes local energetics and friction and does not prove entropy causation.
Conversely, their realistic molecular chemistry and experimental validation
provide application evidence that our coarse-grained study does not offer.
Neither study establishes a universal species-level entropy--diffusion law.

The most informative checks would be:
\begin{enumerate}
\item Reproduce the published local pair statistic with an explicit species,
weighting, density, and radial-range convention, then compare it with a full
all-pair $s_2$ and an independent total-entropy benchmark. For our systems,
calculate the corresponding divalent-ion-centered contributions and compare
both their increments and the remainder $s_{\rm ex}-s_2$ against TI.
\item Repeat the semilog fits with reduced diffusion, independent replicas,
MSD-window checks, uncertainty in both coordinates, and additional compositions
or state points reserved for prediction. A four-formulation fit is a useful
observation, not a strong universality test.
\item Test solvation-class convergence and composition-preserving permutation
references. Compare topology, charge-only classes, and full chemical classes
before interpreting a larger state alphabet as increased thermodynamic entropy.
\item Evaluate association and residence times, viscosity or friction proxies,
and species-resolved motion. Controlled substitutions at matched nominal
ionic strength and concentration can help identify the mechanism, while
retaining a clear account of the interactions that actually change.
\end{enumerate}

\noindent\textbf{Assessment.} This is a relevant and experimentally supported
electrolyte-design study with a useful local structure--mobility correlation.
The main article's evidence is insufficient for its stronger claim of strict
Rosenfeld scaling with total thermodynamic excess entropy or proof of
entropy-driven diffusion. It is a close precedent for our motivation and a
valuable source of local descriptors, rather than a quantitative validation
of the same entropy observable we obtain by HNC and MD/TI.

\section{Voronoi and graph-based structural entropy}
\subsection{Local topology and its implementation}
Lazar, Han, and Srolovitz~\cite{Lazar2015} describe a Voronoi-topological
framework for classifying local atomic structure. The VoroTop
implementation~\cite{LazarToolkit} makes this approach accessible for
atomistic configurations. Graph connectivity resolves distinctions that
coordination numbers or simple face-count indices can miss. These sources
support the structural representation; they do not by themselves validate
an absolute thermodynamic-entropy estimator for our mixtures.

\subsection{Yoon et al.: topological entropy and transport (2019)}
Yoon et al.~\cite{Yoon2019} classify complete Voronoi-cell graph topologies
using Weinberg vectors and calculate a Shannon entropy from their frequencies.
For repulsive soft-sphere fluids, they compare this descriptor with TI excess
entropy and reduced transport. They use an ideal-gas topological reference
and address unseen structural types. The relation to thermodynamic entropy
is interaction dependent, and neighboring-cell correlations remain important.

This is the closest methodological precedent for testing topology against
our TI and diffusion results. It motivates a structural benchmark rather
than assuming that the raw local Shannon entropy equals $s_{\rm ex}/k_B$.
In our mixtures, a further question is whether geometry alone misses changes
in the charge and species identities of neighboring particles. That is an
inference motivating our proposed colored-neighborhood analysis, not a
result already established by this reference.

\subsection{Zhou and Milner: cluster-entropy extrapolation (2016)}
Zhou and Milner~\cite{Zhou2016} classify Voronoi-derived clusters using graph
isomorphism. They estimate subsystem structural entropy and its growth with
particle number; the slope gives an entropy per particle. Their hard-particle
results are consistent with previous configurational-entropy estimates from
thermodynamic integration.

This goes beyond a histogram of isolated cell types by analyzing larger
structures. Its transfer to attractive charged liquids requires separate
validation: graph-equivalent configurations can have different distances,
energies, and within-class phase-space volumes. Sampling rare graph classes
also becomes harder as cluster size grows. The companion plan therefore
makes extrapolation conditional on adequate class coverage, a stable
entropy-growth regime, and an explicit definition of the entropy target.
A stable graph-entropy slope alone is not proof of equality to liquid excess
entropy.

\subsection{Molecular-liquid applications}
Abdulkareem U and Madhurima~\cite{Abdulkareem2024} use Voronoi entropy to
characterize hydrogen-bonded liquids and binary alcohol mixtures. This
establishes broader use of structural-diversity measures in molecular
mixtures. It is not a direct demonstration that they provide absolute excess
entropy or a universal ion-diffusion predictor. The issue year is 2024;
the article became available online in November 2023.

\section{The 2PT comparison and its limits}
Sun et al.~\cite{Sun2017} reanalyze the two-phase thermodynamic approach to
liquid entropy. Lai et al.~\cite{Lai2012} extend 2PT treatment to mixtures.
These are the methodological sources for our corrected classical spectral
assessment, not independent validations of the present charged mixtures.

Our rehabilitation tests pass internal normalization and reference checks,
but retain composition-dependent discrepancies against TI. In particular,
the primary corrected estimate overstates the controlled entropy increment
with monovalent background by approximately $0.036k_B$ per particle.
Consequently TI remains our primary entropy coordinate. A static Voronoi/graph
comparison has the advantage that it does not take a diffusion coefficient
as input, whereas the 2PT gas/solid partition is informed by diffusion.
Agreement between 2PT entropy and diffusion therefore has a different
interpretation from an independently calculated static-structure correlation.

\section{How the prior work connects to our HNC/MD study}
\subsection{The role of HNC and the independent MD test}
Our workflow uses hypernetted-chain (HNC) calculations to screen charge,
composition, and solvent-interaction candidates for higher excess entropy,
then tests selected states with independent MD thermodynamics and transport.
The successful entropy ordering is a validation of the screening for these
candidates, not an exact prediction of their dynamics. In the completed study,
HNC absolute entropies are more negative than TI values by approximately
$0.064$--$0.093k_B$ per particle; controlled diversification increments agree
more closely: about $0.1107$--$0.1108k_B$ from HNC versus
$0.1032$--$0.1033k_B$ from MD/TI.

This makes the thermodynamic and kinetic steps separately assessable.
Agreement of HNC with TI is distinct from agreement of an entropy-based
mobility forecast with measured $D_i$. HNC does not directly provide
self-diffusion coefficients in the present implementation. A species-specific
mobility prediction requires an additional dynamical assumption or model.
The screened candidate set also does not establish a global composition optimum.

\subsection{Connection to the five closest entropy--diffusion precedents}
The links below are our assessment of how the cited studies inform this
project; they are not claims that those authors examined our particular
controls or divalent-solvent model.
\begin{longtable}{@{}p{0.25\textwidth}p{0.69\textwidth}@{}}
\toprule
Prior work & Connection and distinction for the present calculations\\
\midrule
\endhead
Rosenfeld, Nardi, and Zinamon (1995)~\cite{Rosenfeld1995} &
Establishes an early charged-mixture precedent for entropy-based diffusion
modeling. Our neutral solvent, oppositely charged probes, and physical
solvent-diversification controls address a different setting; charged-mixture
scaling itself is not new.\\[3pt]
Krekelberg et al. (2009)~\cite{Krekelberg2009} &
Provides the main context for tracer mobility and component decoupling.
It motivates comparing bulk and probe responses separately. Our ordinary
reduced-diffusion plots do not implement their generalized dilute-limit
scaling, so disagreement with a simple bulk-entropy correlation would not
by itself falsify their model.\\[3pt]
Wang et al. (2023)~\cite{Wang2023} &
Directly motivates the multispecies-electrolyte question: can composition
and solvation diversity improve ion motion? Our simpler model adds a label
null and matched mean bare interaction amplitude. These help examine
attribution, while actual energy, friction, and local solvation still change.\\[3pt]
Schmitt, Hasse, and Stephan (2025)~\cite{Schmitt2025} &
Shows how mixture entropy can support diffusion prediction when supplemented
by limiting transport properties and mixing rules. This supports using HNC
as thermodynamic input while recognizing that it supplies only part of a
predictive transport model. Their binary framework is not directly applied
here to our multicomponent charged mixtures.\\[3pt]
He et al. (2026)~\cite{He2026} &
The closest broad application precedent, now assessed from the supplied
full article (Sec.~\ref{sec:HeAssessment}). Its four-formulation fit uses a
Li-centered pair-entropy proxy; our HNC/MD comparison uses total TI entropy.
Its chemistry and experiments add application evidence, while our independent
entropy benchmark and label control test different methodological questions.\\
\bottomrule
\end{longtable}

\subsection{Would using pair entropy change our conclusions?}
\label{sec:ourPairEntropy}
The existing fine-grid HNC output already contains the \emph{full all-species}
pair entropy. It therefore permits a preliminary comparison without new MD
or a fitted entropy model. Table~\ref{tab:ourPairEntropy} uses these existing
values. This is not an RDF-based MD pair-entropy benchmark: the latter still
needs to be computed independently from the trajectories. In particular,
subtracting HNC $s_2$ from MD/TI $s_{\rm ex}$ would mix two methods and would
not measure an MD higher-order entropy contribution.

\begin{table}[ht]
\centering\small
\caption{Existing HNC pair and full-entropy increments, in $k_B$ per particle.
Differences use the four-decimal fine-grid summary; percentages are descriptive,
not uncertainty estimates.}
\label{tab:ourPairEntropy}
\begin{tabular}{@{}lrrr@{}}
\toprule
Contrast & $\Delta s_2^{\rm HNC}$ & $\Delta s_{\rm ex}^{\rm HNC}$ & Pair fraction\\
\midrule
L5--U4 & 0.0000 & 0.0000 & ---\\
M5--U4 & 0.0226 & 0.0332 & 68\%\\
H5--U4 & 0.0750 & 0.1107 & 68\%\\
H7--U6 & 0.0751 & 0.1108 & 68\%\\
B5--B3 & $-0.0025$ & $-0.0042$ & 60\%\\
\bottomrule
\end{tabular}
\end{table}

For all seven mechanically distinct states, HNC $s_2$ and HNC $s_{\rm ex}$
have the same ordering. The identical-label U4/L5 comparison remains an
exact null. Physical solvent diversification still increases entropy;
monovalent addition B3$\rightarrow$B5 still decreases it slightly. Thus
\emph{the existing HNC pair data support the same qualitative conclusions},
including the broad association with increased diffusion of both the mean
liquid and the divalent probes. The diffusion values and their uncertainties
do not change when the entropy coordinate is replaced, so the small
controlled probe effect remains provisional.

The quantitative interpretation does change. HNC $s_2$ accounts for only
approximately 75.6--76.2\% of the magnitude of HNC $s_{\rm ex}$ across these
states. More relevantly, it captures only about 68\% of the controlled
solvent-diversification increments. The HNC remainder
$s_{\rm ex}^{\rm HNC}-s_2^{\rm HNC}$ increases by about $0.0357k_B$ for
both strong splits. These are differences within the HNC approximation,
not independently measured exact many-particle entropy terms.
Consequently, pair entropy alone would understate the HNC entropy increment
by about one third and change fitted entropy--diffusion slopes. It would
not establish quantitative equality to TI or a universal Rosenfeld relation.
The frequently invoked greater-than-90\% approximation is not supported
by these particular HNC results; this observation does not determine its
accuracy in the different molecular electrolytes of He et al.

Finally, this all-species $s_2$ comparison does not validate the selected
Li-centered statistic used in the main-text description of He et al.
Selected probe-centered contributions can behave differently from the full
pair sum. A useful next check is to calculate both from the same MD RDFs,
with identical density, normalization, and integration conventions, and
compare their controlled increments against TI. Numeric sources and the
preliminary comparisons are recorded in
\path{doc/hnc_pair_entropy_comparison.json}.

\subsection{What this implies for the original question}
An increase in the number of components can increase ideal mixing entropy
without changing forces, as our U4/L5 control demonstrates. To affect
mobility, physical diversification must change the mechanical environment;
that change may simultaneously alter excess entropy and diffusion. Our broad
mixture trends are consistent with an entropy--mobility association, including
for the divalent probes. They do not show that an ideal mixing contribution
causes faster motion, or that one bulk entropy determines each species'
diffusivity. The smaller controlled probe response remains less certain than
the corresponding entropy and solvent-dominated mean-diffusion changes.

\section{Assessment of the present study's contribution}
The following statements concern our completed calculations, not the external
papers. All states have $N=4000$, $T^*=2$, $\rho^*=0.7$, and 3\% of each
divalent sign, with three replicas. The TI reference is the ideal mixture
of the same composition. The LJ Hamiltonian is held fixed along each density
integration path.

\begin{longtable}{@{}p{0.27\textwidth}p{0.67\textwidth}@{}}
\toprule
Finding or control & Interpretation in light of prior work\\
\midrule
\endhead
Identical-label splitting & U4 and L5 have identical interactions. Their
formal ideal mixing entropy differs, while diffusion is identical and TI
entropy agrees within sampling uncertainty. Component labels alone supply
no mechanical acceleration.\\[3pt]
Broad design changes & Higher-excess-entropy solvent designs increase pooled
divalent self-diffusion by about 18--20\% relative to the baselines. This
supports an association but also changes mean interactions and local friction.\\[3pt]
Controlled interaction diversity & At fixed charged composition, nominal
sizes, masses, ionic strength, and mean bare LJ amplitude, strong solvent
splitting raises TI entropy by about $0.103k_B$. Probe diffusion increases
by about 5\% in point estimates, with unresolved replica or fitting-window
uncertainty. Actual energy and friction are not held fixed.\\[3pt]
Species versus bulk response & The controlled increase in weighted mean
self-diffusion is about 8--9\%, larger than the probe response. Broad
entropy--diffusion trends exist for both, but a single bulk entropy does
not determine each species' mobility.\\[3pt]
Independent entropy benchmark & TI permits a direct assessment of HNC,
2PT, and prospective graph descriptors. Passing internal consistency tests
is distinct from agreeing quantitatively with TI across compositions.\\
\bottomrule
\end{longtable}

These controls provide a plausible contribution alongside the existing
literature. The appropriate claim is a controlled examination of the
entropy--mobility association and its limits in model charged mixtures.
The present evidence does not demonstrate causation by entropy, universality
across mixtures, or a converged small enhancement of probe mobility under
every controlled intervention.

For the Voronoi follow-up, the main opportunity is to distinguish local
packing from chemical organization around the probes, and to determine
whether either adds explanatory value beyond TI and pair entropy. Even a
negative benchmark can be useful if a descriptor tracks solvent-dominated
mean diffusion but misses probe behavior. A focused pilot is justified;
cluster extrapolation should follow only after its sampling requirements
and entropy interpretation are demonstrated.

\subsection*{Local sources and reproducibility}
The completed numerical evidence is in the study directory
\path{ComponentEntropyStudy_20261030/}. Relevant files within it are:
\begin{itemize}
\item \path{results/conclusions.md};
\item \path{results/summary.json} and \path{results/diagnostics.json};
\item \path{two_pt_rehabilitation/report.md}.
\end{itemize}
The proposed analysis is documented separately in
\path{doc/voronoi_entropy_work_plan.tex}.
The prior-work assessment does not replace a full reading of the primary
papers or establish that no equivalent controlled study exists elsewhere.

\begin{thebibliography}{99}
\small
\bibitem{Rosenfeld1977}
Y. Rosenfeld, ``Relation between the transport coefficients and the internal
entropy of simple systems,'' \textit{Phys. Rev. A} \textbf{15}, 2545--2549 (1977).
\href{https://doi.org/10.1103/PhysRevA.15.2545}{doi:10.1103/PhysRevA.15.2545}.

\bibitem{Rosenfeld1995}
Y. Rosenfeld, E. Nardi, and Z. Zinamon, ``Corresponding States Hard-Sphere Model
for the Diffusion Coefficients of Binary Dense-Plasma Mixtures,''
\textit{Phys. Rev. Lett.} \textbf{75}, 2490 (1995).
\href{https://doi.org/10.1103/PhysRevLett.75.2490}{doi:10.1103/PhysRevLett.75.2490}.

\bibitem{Krekelberg2009}
W. P. Krekelberg, M. J. Pond, G. Goel, V. K. Shen, J. R. Errington, and
T. M. Truskett, ``Generalized Rosenfeld scalings for tracer diffusivities in
not-so-simple fluids: Mixtures and soft particles,'' \textit{Phys. Rev. E}
\textbf{80}, 061205 (2009).
\href{https://doi.org/10.1103/PhysRevE.80.061205}{doi:10.1103/PhysRevE.80.061205};
\href{https://arxiv.org/abs/0910.0280}{accessible preprint}.

\bibitem{Schmitt2025}
S. Schmitt, H. Hasse, and S. Stephan, ``Entropy scaling for diffusion
coefficients in fluid mixtures,'' \textit{Nature Communications}
\textbf{16}, 2611 (2025).
\href{https://doi.org/10.1038/s41467-025-57780-z}{doi:10.1038/s41467-025-57780-z}.

\bibitem{Wang2023}
Q. Wang, C. Zhao, J. Wang, Z. Yao, S. Wang, S. G. H. Kumar, S. Ganapathy,
S. Eustace, X. Bai, B. Li, and M. Wagemaker,
``High entropy liquid electrolytes for lithium batteries,''
\textit{Nature Communications} \textbf{14}, 440 (2023).
\href{https://doi.org/10.1038/s41467-023-36075-1}{doi:10.1038/s41467-023-36075-1}.

\bibitem{He2026}
L. He, Y. Lin, J. Li, H. Su, H. Liu, C. Yang, and C. Lian,
``Bridging thermodynamics and ion transport in lithium batteries via molecular
design of high-entropy electrolytes,'' \textit{Molecular Systems Design \&
Engineering} \textbf{11}, 681--690 (2026).
\href{https://doi.org/10.1039/D6ME00060F}{doi:10.1039/D6ME00060F}.
First published 19 May 2026. Full main article assessed from the supplied
\path{He2026.pdf}; supplementary information not assessed.

\bibitem{Yoon2019}
T. J. Yoon, M. Y. Ha, E. A. Lazar, W. B. Lee, and Y.-W. Lee,
``Topological extension of the isomorph theory based on the Shannon entropy,''
\textit{Phys. Rev. E} \textbf{100}, 012118 (2019).
\href{https://doi.org/10.1103/PhysRevE.100.012118}{doi:10.1103/PhysRevE.100.012118};
\href{https://arxiv.org/abs/1901.02772}{accessible preprint}.

\bibitem{Lazar2015}
E. A. Lazar, J. Han, and D. J. Srolovitz, ``A topological framework for local
structure analysis in condensed matter,'' \textit{Proceedings of the National
Academy of Sciences USA} \textbf{112}, E5769--E5776 (2015).
\href{https://doi.org/10.1073/pnas.1505788112}{doi:10.1073/pnas.1505788112};
\href{https://arxiv.org/abs/1508.05937}{accessible preprint}.

\bibitem{LazarToolkit}
E. A. Lazar, ``VoroTop: Voronoi Cell Topology Visualization and Analysis
Toolkit,'' \textit{Modelling and Simulation in Materials Science and
Engineering} \textbf{26}, 015011 (2018; published online December 2017).
\href{https://doi.org/10.1088/1361-651X/aa9a01}{doi:10.1088/1361-651X/aa9a01};
\href{https://arxiv.org/abs/1804.04221}{accessible preprint};
\href{https://www.vorotop.org/}{software and documentation}.

\bibitem{Zhou2016}
Y. Zhou and S. T. Milner, ``Structural entropy of glassy systems from graph
isomorphism,'' \textit{Soft Matter} \textbf{12}, 7281--7288 (2016).
\href{https://doi.org/10.1039/C6SM01355D}{doi:10.1039/C6SM01355D}.
Primary source for cluster-entropy extrapolation.

\bibitem{Abdulkareem2024}
Abdulkareem U and V. Madhurima, ``Understanding complexity of hydrogen-bonded
liquids through Voronoi entropy,'' \textit{Chemical Physics Impact}
\textbf{8}, 100403 (2024; available online 30 November 2023).
\href{https://doi.org/10.1016/j.chphi.2023.100403}{doi:10.1016/j.chphi.2023.100403}.

\bibitem{Sun2017}
T. Sun, J. Xian, H. Zhang, Z. Zhang, and Y. Zhang,
``Two-phase thermodynamic model for computing entropies of liquids reanalyzed,''
\textit{J. Chem. Phys.} \textbf{147}, 194505 (2017).
\href{https://doi.org/10.1063/1.5001798}{doi:10.1063/1.5001798};
\href{https://arxiv.org/abs/1711.02041}{accessible preprint}.

\bibitem{Lai2012}
P.-K. Lai, C.-M. Hsieh, and S.-T. Lin,
``Rapid determination of entropy and free energy of mixtures from molecular
dynamics simulations with the two-phase thermodynamic model,''
\textit{Phys. Chem. Chem. Phys.} \textbf{14}, 15206--15213 (2012).
\href{https://doi.org/10.1039/C2CP42011B}{doi:10.1039/C2CP42011B}.

\bibitem{Dyre2018}
J. C. Dyre, ``Perspective: Excess-entropy scaling,''
\textit{J. Chem. Phys.} \textbf{149}, 210901 (2018).
\href{https://doi.org/10.1063/1.5055064}{doi:10.1063/1.5055064};
\href{https://jcdyre.dk/2018_JCP_149_210901.pdf}{author-hosted full text}.

\bibitem{WangAdvanced2023}
Q. Wang, C. Zhao, Z. Yao, J. Wang, F. Wu, S. G. H. Kumar, S. Ganapathy,
S. Eustace, X. Bai, B. Li, J. Lu, and M. Wagemaker,
``Entropy-Driven Liquid Electrolytes for Lithium Batteries,''
\textit{Advanced Materials} \textbf{35}, 2210677 (2023).
\href{https://doi.org/10.1002/adma.202210677}{doi:10.1002/adma.202210677}.
\end{thebibliography}
\end{document}
